\section{Sheafification}
\label{section-sheafification}

\noindent
In order to define the sheafification we study the zeroth
{\v C}ech cohomology group of a covering and its functoriality
properties.

\medskip\noindent
Let $\mathcal{F}$ be a presheaf of sets on $\mathcal{C}$, and let
$\mathcal{U} = \{U_i \to U\}_{i \in I}$ be a covering of $\mathcal{C}$.
Let us use the notation $\mathcal{F}(\mathcal{U})$ to indicate the equalizer
$$
H^0(\mathcal{U}, \mathcal{F})
=
\{
(s_i)_{i\in I} \in \prod\nolimits_i \mathcal{F}(U_i)
\mid
s_i|_{U_i \times_U U_j} = s_j|_{U_i \times_U U_j}
\ \forall i, j \in I
\}.
$$
As we will see later, this is the zeroth {\v C}ech cohomology
of $\mathcal{F}$ over $U$ with respect to the covering $\mathcal{U}$.
A small remark is that we can define $H^0(\mathcal{U}, \mathcal{F})$
as soon as all the morphisms $U_i \to U$ are representable, i.e.,
$\mathcal{U}$ need not be a covering of the site.
There is a canonical map $\mathcal{F}(U) \to H^0(\mathcal{U}, \mathcal{F})$.
It is clear that a morphism of coverings $\mathcal{U} \to \mathcal{V}$
induces commutative diagrams
$$
\xymatrix{
& U_i \ar[rr] & & V_{\alpha(i)} \\
U_i \times_U U_j \ar[rr] \ar[ur] \ar[dr] & &
V_{\alpha(i)} \times_V V_{\alpha(j)} \ar[ur] \ar[dr] & \\
& U_j \ar[rr] & & V_{\alpha(j)}
}.
$$
This in turn produces a map $H^0(\mathcal{V}, \mathcal{F}) \to
H^0(\mathcal{U}, \mathcal{F})$, compatible with the map $\mathcal{F}(V)
\to \mathcal{F}(U)$.

\medskip\noindent
By construction, a presheaf $\mathcal{F}$ is a sheaf if and only if for
every covering $\mathcal{U}$ of $\mathcal{C}$ the natural map
$\mathcal{F}(U) \to H^0(\mathcal{U}, \mathcal{F})$ is bijective.
We will use this notion to prove the following
simple lemma about limits of sheaves.

\begin{lemma}
\label{lemma-limit-sheaf}
\begin{slogan}
Limit of diagram of sheaves exists and coincides with limit as presheaves
\end{slogan}
Let $\mathcal{F} : \mathcal{I} \to \Sh(\mathcal{C})$
be a diagram. Then $\lim_\mathcal{I} \mathcal{F}$ exists
and is equal to the limit in the category of presheaves.
\end{lemma}

\begin{proof}
Let $\lim_i \mathcal{F}_i$ be the limit as a presheaf.
We will show that this is a sheaf and then it will trivially follow
that it is a limit in the category of sheaves. To prove the sheaf
property, let $\mathcal{V} = \{V_j \to V\}_{j\in J}$ be a covering.
Let $(s_j)_{j\in J}$ be an element of $H^0(\mathcal{V}, \lim_i \mathcal{F}_i)$.
Using the projection maps we get elements $(s_{j, i})_{j\in J}$
in $H^0(\mathcal{V}, \mathcal{F}_i)$. By the sheaf property for
$\mathcal{F}_i$ we see that there is a unique $s_i \in \mathcal{F}_i(V)$
such that $s_{j, i} = s_i|_{V_j}$. Let $\phi : i \to i'$ be a morphism
of the index category. We would like to show that
$\mathcal{F}(\phi) : \mathcal{F}_i \to \mathcal{F}_{i'}$
maps $s_i$ to $s_{i'}$. We know this is true for the sections
$s_{i, j}$ and $s_{i', j}$ for all $j$ and hence by the sheaf property
for $\mathcal{F}_{i'}$ this is true. At this point we have an
element $s = (s_i)_{i \in \Ob(\mathcal{I})}$ of
$(\lim_i \mathcal{F}_i)(V)$. We leave it to the reader to see
this element has the required property that $s_j = s|_{V_j}$.
\end{proof}

\begin{example}
\label{example-singleton-sheaf}
A particular example is the limit over the empty diagram.
This gives the final object in the category of (pre)sheaves.
It is the presheaf that associates
to each object $U$ of $\mathcal{C}$ a singleton set, with unique
restriction mappings and moreover this presheaf is a sheaf.
We often denote this sheaf by $*$.
\end{example}

\noindent
Let $\mathcal{J}_U$ be the category of all coverings of $U$.
In other words, the objects of $\mathcal{J}_U$ are the coverings
of $U$ in $\mathcal{C}$, and the morphisms are the refinements.
By our conventions on sites this is indeed a category, i.e.,
the collection of objects and morphisms forms a set.
Note that $\Ob(\mathcal{J}_U)$ is not empty since
$\{\text{id}_U\}$ is an object of it. According to the remarks
above the construction $\mathcal{U} \mapsto H^0(\mathcal{U}, \mathcal{F})$
is a contravariant functor on $\mathcal{J}_U$.
We define
$$
\mathcal{F}^{+}(U)
=
\colim_{\mathcal{J}_U^{opp}}
H^0(\mathcal{U}, \mathcal{F})
$$
See Categories, Section \ref{categories-section-limits} for
a discussion of limits and colimits. We point out that later
we will see that $\mathcal{F}^{+}(U)$ is the zeroth {\v C}ech
cohomology of $\mathcal{F}$ over $U$.

\medskip\noindent
Before we say more about the structure of the colimit, we turn
the collection of sets
$\mathcal{F}^{+}(U)$, $U \in \Ob(\mathcal{C})$
into a presheaf. Namely, let $V \to U$ be a morphism of $\mathcal{C}$.
By the axioms of a site there is a functor\footnote{This construction
actually involves a choice of the fibre products $U_i \times_U V$
and hence the axiom of choice. The resulting map does not depend on
the choices made, see below.}
$$
\mathcal{J}_U
\longrightarrow
\mathcal{J}_V, \quad
\{U_i \to U\}
\longmapsto
\{U_i \times_U V \to V\}.
$$
Note that the projection maps furnish a functorial
morphism of coverings $\{U_i \times_U V \to V\} \to \{U_i \to U\}$
and hence, by the construction above, a functorial map of sets
$H^0(\{U_i \to U\}, \mathcal{F}) \to
H^0(\{U_i \times_U V \to V\}, \mathcal{F})$.
In other words, there is a transformation of functors
from $H^0(-, \mathcal{F}) : \mathcal{J}_U^{opp} \to \textit{Sets}$
to the composition
$\mathcal{J}_U^{opp} \to \mathcal{J}_V^{opp}
\xrightarrow{H^0(-, \mathcal{F})} \textit{Sets}$. Hence by
generalities of colimits we obtain a canonical map
$\mathcal{F}^+(U) \to \mathcal{F}^+(V)$. In terms of the description
of the set $\mathcal{F}^+(U)$ above, it just takes the element
associated with $s = (s_i) \in H^0(\{U_i \to U\}, \mathcal{F})$ to the
element associated with $(s_i|_{V \times_U U_i})
\in H^0(\{U_i \times_U V \to V\}, \mathcal{F})$.

\begin{lemma}
\label{lemma-plus-presheaf}
The constructions above define a presheaf
$\mathcal{F}^+$ together with a canonical
map of presheaves $\mathcal{F} \to \mathcal{F}^+$.
\end{lemma}

\begin{proof}
All we have to do is to show that given morphisms
$W \to V \to U$ the composition $\mathcal{F}^+(U)
\to \mathcal{F}^+(V) \to \mathcal{F}^+(W)$
equals the map $\mathcal{F}^+(U) \to \mathcal{F}^+(W)$.
This can be shown directly by verifying that, given
a covering $\{U_i \to U\}$ and
$s = (s_i) \in H^0(\{U_i \to U\}, \mathcal{F})$,
we have canonically
$W \times_U U_i \cong W \times_V (V \times_U U_i)$,
and
$s_i|_{W \times_U U_i}$
corresponds to
$(s_i|_{V \times_U U_i})|_{W \times_V (V \times_U U_i)}$
via this isomorphism.
\end{proof}

\noindent
More indirectly, the result of
Lemma \ref{lemma-independent-refinement} shows that
we may pullback an element $s$ as above via any morphism
from any covering of $W$ to $\{U_i \to U\}$
and we will always end up with the same element in
$\mathcal{F}^+(W)$.

\begin{lemma}
\label{lemma-plus-functorial}
The association $\mathcal{F} \mapsto
(\mathcal{F} \to \mathcal{F}^+)$
is a functor.
\end{lemma}

\begin{proof}
Instead of proving this we state exactly what needs to be proven.
Let $\mathcal{F} \to \mathcal{G}$ be a map of presheaves. Prove
the commutativity of:
$$
\xymatrix{
\mathcal{F} \ar[r] \ar[d]
&
\mathcal{F}^{+} \ar[d]
\\
\mathcal{G} \ar[r]
&
\mathcal{G}^{+}
}
$$
\end{proof}

\noindent
The next two lemmas imply that the colimits above are colimits
over a directed set.

\begin{lemma}
\label{lemma-common-refinement}
Given a pair of coverings $\{U_i \to U\}$
and $\{V_j \to U\}$ of a given object $U$ of the site
$\mathcal{C}$, there exists a covering which is a
common refinement.
\end{lemma}

\begin{proof}
Since $\mathcal{C}$ is a site we have that for every $i$ the
family $\{V_j \times_U U_i \to U_i\}_j$ is a covering.
And, then another axiom implies that $\{V_j \times_U U_i \to U\}_{i, j}$
is a covering of $U$. Clearly this covering refines both given
coverings.
\end{proof}

\begin{lemma}
\label{lemma-independent-refinement}
Any two morphisms $f, g: \mathcal{U} \to \mathcal{V}$ of coverings
inducing the same morphism $U \to V$ induce the same
map $H^0(\mathcal{V}, \mathcal{F}) \to  H^0(\mathcal{U}, \mathcal{F})$.
\end{lemma}

\begin{proof}
Let $\mathcal{U} = \{U_i \to U\}_{i\in I}$ and
$\mathcal{V} = \{V_j \to V\}_{j\in J}$.
The morphism $f$ consists of a map $U\to V$, a map $\alpha : I \to J$ and
maps $f_i : U_i \to V_{\alpha(i)}$.
Likewise, $g$~determines a map $\beta : I \to J$ and maps
$g_i : U_i \to V_{\beta(i)}$.
As $f$ and $g$ induce the same map $U\to V$, the diagram
$$
\xymatrix{
&
V_{\alpha(i)} \ar[dr]
\\
U_i \ar[ur]^{f_i} \ar[dr]_{g_i}
&
&
V
\\
&
V_{\beta(i)} \ar[ur]
}
$$
is commutative for every $i\in I$. Hence $f$ and $g$ factor through
the fibre product
$$
\xymatrix{
&
V_{\alpha(i)}
\\
U_i \ar[r]^-\varphi \ar[ur]^{f_i} \ar[dr]_{g_i}
&
V_{\alpha(i)} \times_V V_{\beta(i)} \ar[u]_{\text{pr}_1} \ar[d]^{\text{pr}_2}
\\
&
V_{\beta(i)}.
}
$$
Now let $s = (s_j)_j \in H^0(\mathcal{V}, \mathcal{F})$.
Then for all $i\in I$:
$$
(f^*s)_i
=
f_i^*(s_{\alpha(i)})
=
\varphi^*\text{pr}_1^*(s_{\alpha(i)})
=
\varphi^*\text{pr}_2^*(s_{\beta(i)})
=
g_i^*(s_{\beta(i)})
=
(g^*s)_i,
$$
where the middle equality is given by the definition
of $H^0(\mathcal{V}, \mathcal{F})$.
This shows that the maps
$H^0(\mathcal{V}, \mathcal{F}) \to H^0(\mathcal{U}, \mathcal{F})$
induced by $f$ and $g$ are equal.
\end{proof}

\begin{remark}
\label{remark-both-refine-same-H0}
In particular this lemma shows that if $\mathcal{U}$ is
a refinement of $\mathcal{V}$, and if $\mathcal{V}$ is a
refinement of $\mathcal{U}$, then there is a canonical
identification $H^0(\mathcal{U}, \mathcal{F}) =
H^0(\mathcal{V}, \mathcal{F})$.
\end{remark}

\noindent
From these two lemmas, and the fact that $\mathcal{J}_U$ is nonempty,
it follows that the diagram $H^0(-, \mathcal{F}) : \mathcal{J}_U^{opp}
\to \textit{Sets}$ is filtered, see
Categories, Definition \ref{categories-definition-directed}.
Hence, by Categories,
Section \ref{categories-section-directed-colimits}
the colimit $\mathcal{F}^{+}(U)$ may be described
in the following straightforward manner. Namely, every element in the set
$\mathcal{F}^{+}(U)$ arises from an element
$s \in H^0(\mathcal{U}, \mathcal{F})$ for some covering
$\mathcal{U}$ of $U$. Given a second element $s' \in
H^0(\mathcal{U}', \mathcal{F})$ then $s$ and $s'$ determine
the same element of the colimit if and only if there exists a covering
$\mathcal{V}$ of $U$ and refinements $f : \mathcal{V} \to \mathcal{U}$ and
$f' : \mathcal{V} \to \mathcal{U}'$ such that $f^*s = (f')^*s'$
in $H^0(\mathcal{V}, \mathcal{F})$. Since the trivial covering
$\{\text{id}_U\}$ is an object of $\mathcal{J}_U$ we get
a canonical map $\mathcal{F}(U) \to \mathcal{F}^+(U)$.

\begin{lemma}
\label{lemma-plus-surjective}
The map $\theta : \mathcal{F} \to \mathcal{F}^+$ has the following
property: For every object $U$ of $\mathcal{C}$ and every section
$s \in \mathcal{F}^+(U)$ there exists a covering $\{U_i \to U\}$
such that $s|_{U_i}$ is in the image of $\theta : \mathcal{F}(U_i)
\to \mathcal{F}^{+}(U_i)$.
\end{lemma}

\begin{proof}
Namely, let $\{U_i \to U\}$ be a covering such that $s$ arises
from the element $(s_i) \in H^0(\{U_i \to U\}, \mathcal{F})$.
According to Lemma \ref{lemma-independent-refinement} we may
consider the covering $\{U_i \to U_i\}$ and the (obvious) morphism
of coverings $\{U_i \to U_i\} \to \{U_i \to U\}$ to compute the
pullback of $s$ to an element of $\mathcal{F}^+(U_i)$. And indeed,
using this covering we get exactly $\theta(s_i)$ for the restriction
of $s$ to $U_i$.
\end{proof}

\begin{definition}
\label{definition-separated}
We say that a presheaf of sets $\mathcal{F}$ on a site
$\mathcal{C}$ is {\it separated} if, for all coverings
of $\{U_i \rightarrow U\}$, the map
$\mathcal{F}(U) \to \prod \mathcal{F}(U_i)$ is injective.
\end{definition}

\begin{theorem}
\label{theorem-plus}
With $\mathcal{F}$ as above
\begin{enumerate}
\item
\label{item-sep}
The presheaf $\mathcal{F}^+$ is separated.
\item
\label{item-sheaf}
If $\mathcal{F}$ is separated, then $\mathcal{F}^+$ is a sheaf
and the map of presheaves $\mathcal{F} \to \mathcal{F}^+$ is injective.
\item
\label{item-plus-iso}
If $\mathcal{F}$ is a sheaf, then $\mathcal{F} \to \mathcal{F}^+$
is an isomorphism.
\item
\label{item-plusplus}
The presheaf $\mathcal{F}^{++}$ is always a sheaf.
\end{enumerate}
\end{theorem}

\begin{proof}
Proof of (\ref{item-sep}).
Suppose that $s, s' \in \mathcal{F}^+(U)$ and suppose that
there exists some covering $\{U_i \to U\}$ such that
$s|_{U_i} = s'|_{U_i}$ for all $i$. We now have three coverings
of $U$: the covering $\{U_i \to U\}$ above, a covering $\mathcal{U}$
for $s$ as in Lemma \ref{lemma-plus-surjective},
and a similar covering $\mathcal{U}'$ for $s'$. By Lemma
\ref{lemma-common-refinement}, we can find a common refinement,
say $\{W_j \to U\}$. This means we have $s_j, s'_j \in \mathcal{F}(W_j)$
such that $s|_{W_j} = \theta(s_j)$, similarly for $s'|_{W_j}$, and
such that $\theta(s_j) = \theta(s'_j)$. This last equality means
that there exists some covering $\{W_{jk} \to W_j\}$ such that
$s_j|_{W_{jk}} = s'_j|_{W_{jk}}$. Then since $\{W_{jk} \to U\}$
is a covering we see that $s, s'$ map to the same element of
$H^0(\{W_{jk} \to U\}, \mathcal{F})$ as desired.

\medskip\noindent
Proof of (\ref{item-sheaf}). It is clear that $\mathcal{F} \to
\mathcal{F}^+$ is injective because all the maps
$\mathcal{F}(U) \to H^0(\mathcal{U}, \mathcal{F})$
are injective. It is also clear that, if $\mathcal{U} \to
\mathcal{U}'$ is a refinement, then $H^0(\mathcal{U}', \mathcal{F})
\to H^0(\mathcal{U}, \mathcal{F})$ is injective. Now,
suppose that $\{U_i \to U\}$ is a covering, and let
$(s_i)$ be a family of elements of $\mathcal{F}^+(U_i)$
satisfying the sheaf condition
$s_i|_{U_i \times_U U_{i'}} = s_{i'}|_{U_i \times_U U_{i'}}$
for all $i, i' \in I$. Choose coverings (as in
Lemma \ref{lemma-plus-surjective}) $\{U_{ij} \to U_i\}$
such that $s_i|_{U_{ij}}$ is the image of the (unique)
element $s_{ij} \in \mathcal{F}(U_{ij})$. The sheaf condition
implies that $s_{ij}$ and $s_{i'j'}$ agree over
$U_{ij} \times_U U_{i'j'}$ because it maps to
$U_i \times_U U_{i'}$ and we have the equality there.
Hence $(s_{ij}) \in H^0(\{U_{ij} \to U\}, \mathcal{F})$
gives rise to an element $s \in \mathcal{F}^+(U)$. We leave
it to the reader to verify that $s|_{U_i} = s_i$.

\medskip\noindent
Proof of (\ref{item-plus-iso}). This is immediate from the definitions
because the sheaf property says exactly that every map
$\mathcal{F} \to H^0(\mathcal{U}, \mathcal{F})$ is bijective
(for every covering $\mathcal{U}$ of $U$).

\medskip\noindent
Statement (\ref{item-plusplus}) is now obvious.
\end{proof}

\begin{definition}
\label{definition-associated-sheaf}
Let $\mathcal{C}$ be a site and let $\mathcal{F}$ be a presheaf
of sets on $\mathcal{C}$. The sheaf $\mathcal{F}^\# := \mathcal{F}^{++}$
together with the canonical map $\mathcal{F} \to \mathcal{F}^\#$
is called the {\it sheaf associated to $\mathcal{F}$}.
\end{definition}

\begin{proposition}
\label{proposition-sheafification-adjoint}
The canonical map $\mathcal{F} \to \mathcal{F}^\#$ has the
following universal property: For any map $\mathcal{F} \to \mathcal{G}$,
where $\mathcal{G}$ is a sheaf of sets, there is a unique map
$\mathcal{F}^\# \to \mathcal{G}$ such that $\mathcal{F} \to \mathcal{F}^\#
\to \mathcal{G}$ equals the given map.
\end{proposition}

\begin{proof}
By Lemma \ref{lemma-plus-functorial} we get a commutative diagram
$$
\xymatrix{
\mathcal{F} \ar[r] \ar[d]
&
\mathcal{F}^{+} \ar[r] \ar[d]
&
\mathcal{F}^{++} \ar[d]
\\
\mathcal{G} \ar[r]
&
\mathcal{G}^{+} \ar[r]
&
\mathcal{G}^{++}
}
$$
and by Theorem \ref{theorem-plus} the lower horizontal maps
are isomorphisms. The uniqueness follows from Lemma
\ref{lemma-plus-surjective} which says that every section of
$\mathcal{F}^\#$ locally comes from sections of $\mathcal{F}$.
\end{proof}

\noindent
It is clear from this result that the functor $\mathcal{F}
\mapsto (\mathcal{F} \to \mathcal{F}^\#)$ is unique
up to unique isomorphism of functors. Actually, let us temporarily
denote $i : \Sh(\mathcal{C}) \to \textit{PSh}(\mathcal{C})$
the functor of inclusion. The result above actually says that
$$
\Mor_{\textit{PSh}(\mathcal{C})}(\mathcal{F}, i(\mathcal{G}))
=
\Mor_{\Sh(\mathcal{C})}(\mathcal{F}^\#, \mathcal{G}).
$$
In other words, the functor of sheafification is the left adjoint
to the inclusion functor $i$.

\begin{example}
\label{example-constant-sheaf}
Let $\mathcal{C}$ be a site. Let $E$ be a set.
The {\it constant sheaf $\underline{E}$ with value $E$}
is the sheafification of the presheaf given by the constant
functor with value $E$. It is characterized by the rule
$\Mor_{\Sh(\mathcal{C})}(\underline{E}, \mathcal{F}) =
\Mor_{\textit{Sets}}(E, \Gamma(\mathcal{C}, \mathcal{F}))$
for every sheaf $\mathcal{F}$ on $\mathcal{C}$ where the
notation is from Definition \ref{definition-global-sections}.
\end{example}

\noindent
We finish this section with a couple of lemmas.

\begin{lemma}
\label{lemma-colimit-sheaves}
\begin{slogan}
Colimit in category of sheaves equals sheafification of colimit in
category of presheaves.
\end{slogan}
Let $\mathcal{F} : \mathcal{I} \to \Sh(\mathcal{C})$
be a diagram. Then $\colim_\mathcal{I} \mathcal{F}$ exists
and is the sheafification of the colimit in the category of presheaves.
\end{lemma}

\begin{proof}
Since the sheafification functor is a left adjoint it commutes
with all colimits, see Categories,
Lemma \ref{categories-lemma-adjoint-exact}.
Hence, since $\textit{PSh}(\mathcal{C})$ has colimits, we deduce
that $\Sh(\mathcal{C})$ has colimits (which are the
sheafifications of the colimits in presheaves).
\end{proof}

\begin{lemma}
\label{lemma-sheafification-exact}
The functor $\textit{PSh}(\mathcal{C}) \to \Sh(\mathcal{C})$,
$\mathcal{F} \mapsto \mathcal{F}^\#$ is exact.
\end{lemma}

\begin{proof}
Since it is a left adjoint it is right exact, see
Categories, Lemma \ref{categories-lemma-exact-adjoint}.
On the other hand, by Lemmas \ref{lemma-common-refinement}
and Lemma \ref{lemma-independent-refinement} the colimits
in the construction of $\mathcal{F}^+$ are really over the
directed set $\Ob(\mathcal{J}_U)$ where
$\mathcal{U} \geq \mathcal{U}'$ if and only if
$\mathcal{U}$ is a refinement of $\mathcal{U}'$. Hence by
Categories, Lemma \ref{categories-lemma-directed-commutes}
we see that $\mathcal{F} \to \mathcal{F}^+$ commutes
with finite limits (as a functor from presheaves to
presheaves). Then we conclude using
Lemma \ref{lemma-limit-sheaf}.
\end{proof}

\begin{lemma}
\label{lemma-sections-sheafification}
Let $\mathcal{C}$ be a site.
Let $\mathcal{F}$ be a presheaf of sets on $\mathcal{C}$.
Denote $\theta^2 : \mathcal{F} \to \mathcal{F}^\#$ the canonical
map of $\mathcal{F}$ into its sheafification.
Let $U$ be an object of $\mathcal{C}$.
Let $s \in \mathcal{F}^\#(U)$. There exists
a covering $\{U_i \to U\}$ and sections
$s_i \in \mathcal{F}(U_i)$ such that
\begin{enumerate}
\item $s|_{U_i} = \theta^2(s_i)$, and
\item for every $i, j$ there exists a covering
$\{U_{ijk} \to U_i \times_U U_j\}$ of $\mathcal{C}$ such that
the pullbacks of $s_i$ and $s_j$ to each $U_{ijk}$ agree.
\end{enumerate}
Conversely, given any covering $\{U_i \to U\}$, elements
$s_i \in \mathcal{F}(U_i)$ such that (2) holds, then there
exists a unique section $s \in \mathcal{F}^\#(U)$ such
that (1) holds.
\end{lemma}

\begin{proof}
Omitted.
\end{proof}

\begin{lemma}
\label{lemma-isomorphism-sheafifications}
Let $\mathcal{C}$ be a site. Let $\mathcal{F} \to \mathcal{G}$ be a map
of presheaves of sets on $\mathcal{C}$. Denote $\mathcal{B}$ the set of
$U \in \Ob(\mathcal{C})$ such that $\mathcal{F}(U) \to \mathcal{G}(U)$
is bijective. If every object of $\mathcal{C}$ has a covering by elements
of $\mathcal{B}$, then $\mathcal{F}^\# \to \mathcal{G}^\#$ is an isomorphism.
\end{lemma}

\begin{proof}
Let $U \in \Ob(\mathcal{C})$.
Let us prove that $\mathcal{F}^\#(U) \to \mathcal{G}^\#(U)$ is
surjective. To do this we will use Lemma \ref{lemma-sections-sheafification}
without further mention. For any $s \in \mathcal{G}^\#(U)$ there exists
a covering $\{U_i \to U\}$ and sections $s_i \in \mathcal{G}(U_i)$ such that
\begin{enumerate}
\item $s|_{U_i}$ is the image of $s_i$ via $\mathcal{G} \to \mathcal{G}^\#$,
and
\item for every $i, j$ there exists a covering
$\{U_{ijk} \to U_i \times_U U_j\}$ of $\mathcal{D}$ such that
the pullbacks of $s_i$ and $s_j$ to each $U_{ijk}$ agree.
\end{enumerate}
By assumption, for each $i$ we may choose a covering $\{U_{i, a} \to U_i\}$
with $U_{i, a} \in \mathcal{B}$. Then $\{U_{i, a} \to U\}$ is a covering.
Denoting $s_{i, a}$ the image of $s_i$ in $\mathcal{G}(U_{i, a})$
we see that the pullbacks of $s_{i, a}$ and $s_{j, b}$ to the members
of the covering
$$
\{(U_{i, a} \times_U U_{j, b}) \times_{U_i \times_U U_j} U_{ijk}
\to U_{i, a} \times_U U_{j, b}\}
$$
agree. Hence we may assume that $U_i \in \mathcal{B}$. Repeating
the argument, we may also assume $U_{ijk} \in \mathcal{B}$ for
all $i, j, k$ (details omitted).
Then since $\mathcal{F}(U_i) \to \mathcal{G}(U_i)$
is bijective by our definition of $\mathcal{B}$ in the statement
of the lemma, we get unique $t_i \in \mathcal{F}(U_i)$
mapping to $s_i$. The pullbacks of $t_i$ and $t_j$
to each $U_{ijk}$ agree in $\mathcal{F}(U_{ijk})$ because
$U_{ijk} \in \mathcal{B}$ and because we have the agreement
for $t_i$ and $t_j$.
Then the lemma tells us there exists a unique
$t \in \mathcal{F}^\#(U)$ such that $t|_{U_i}$ is the image of $t_i$.
It follows that $t$ maps to $s$.
We omit the proof that $\mathcal{F}^\#(U) \to \mathcal{G}^\#(U)$
is injective.
\end{proof}








